\documentclass[11pt,a4paper]{article}
% Même mise en forme que le polycopié de prérequis (boîtes, exercices, corrigés).
% ===========================================================================
%  Polycopié de prérequis — Comptabilité prudentielle (M2 alternance)
%  Préambule partagé
%  F. Borel-Mathurin — Université Paris 1 Panthéon-Sorbonne
% ===========================================================================

\usepackage[utf8]{inputenc}
\usepackage[T1]{fontenc}
\IfFileExists{french.ldf}{\usepackage[french]{babel}}{}
\IfFileExists{lmodern.sty}{\usepackage{lmodern}}{}
\usepackage[a4paper,margin=2.4cm]{geometry}

\usepackage{amsmath,amssymb}
\usepackage{booktabs}
\usepackage{array}
\usepackage{enumitem}
\usepackage{xcolor}
\usepackage{tcolorbox}
\tcbuselibrary{skins,breakable}
\usepackage{tikz}
\usetikzlibrary{arrows.meta,positioning,calc}
\usepackage{fancyhdr}
\usepackage[hidelinks]{hyperref}
\usepackage{makeidx}
\makeindex
\IfFileExists{eurosym.sty}{\usepackage{eurosym}}{\providecommand{\euro}{\texteuro}}

% ---------------------------------------------------------------- couleurs
% À remplacer par les valeurs exactes de la charte si besoin.
\definecolor{bleuactu}{RGB}{0,84,150}
\definecolor{vertactu}{RGB}{0,132,100}
\definecolor{rougeactu}{RGB}{175,30,45}
\definecolor{orangeactu}{RGB}{230,140,0}
\definecolor{grisactu}{RGB}{110,110,110}

% ------------------------------------------------- version étudiant/corrigé
% \correctiontrue  -> corrigés inclus
% \correctionfalse -> version étudiant
\newif\ifcorrection
\correctiontrue

% ------------------------------------------------------------------ boîtes
\newtcolorbox{notion}[1]{%
  breakable, enhanced, colback=bleuactu!4, colframe=bleuactu,
  boxrule=0.4pt, left=8pt, right=8pt, top=6pt, bottom=6pt,
  fonttitle=\bfseries\sffamily, coltitle=white, title={#1}}

\newtcolorbox{generalisation}[1]{%
  breakable, enhanced, colback=orangeactu!5, colframe=orangeactu!85!black,
  boxrule=0.4pt, left=8pt, right=8pt, top=6pt, bottom=6pt,
  fonttitle=\bfseries\sffamily, coltitle=white, title={#1}}

\newtcolorbox{pont}[1]{%
  breakable, enhanced, colback=rougeactu!4, colframe=rougeactu,
  boxrule=0.4pt, left=8pt, right=8pt, top=6pt, bottom=6pt,
  fonttitle=\bfseries\sffamily, coltitle=white, title={#1}}

\newtcolorbox{refreg}{%
  breakable, enhanced, colback=grisactu!6, colframe=grisactu!60,
  boxrule=0.4pt, left=8pt, right=8pt, top=6pt, bottom=6pt,
  fontupper=\small}

% --------------------------------------------------------------- exercices
\newcounter{exo}[section]
\renewcommand{\theexo}{\thesection.\arabic{exo}}

\newenvironment{exercice}[1][]{%
  \refstepcounter{exo}%
  \par\medskip\noindent
  \begin{tcolorbox}[breakable, enhanced, colback=white, colframe=vertactu,
     boxrule=0.4pt, left=8pt, right=8pt, top=6pt, bottom=6pt,
     fonttitle=\bfseries\sffamily, coltitle=white,
     title={Exercice \theexo\ifx\relax#1\relax\else\ --- #1\fi}]}%
  {\end{tcolorbox}\par\medskip}

\newenvironment{corrige}{%
  \ifcorrection
    \par\smallskip\noindent
    \begin{tcolorbox}[breakable, enhanced, colback=vertactu!4,
       colframe=vertactu!50, boxrule=0.4pt, left=8pt, right=8pt,
       top=6pt, bottom=6pt, fontupper=\small,
       title={\itshape Corrigé}, fonttitle=\small\itshape,
       coltitle=vertactu!80!black, colbacktitle=vertactu!10]%
  \else
    \setbox0\vbox\bgroup
  \fi}%
  {\ifcorrection\end{tcolorbox}\par\smallskip\else\egroup\fi}

% ------------------------------------------------------------- échéanciers
\tikzset{
  timeaxis/.style={-{Latex[length=2mm]}, thick, grisactu},
  timenode/.style={circle, fill=bleuactu, inner sep=1.1pt},
  flux/.style={-{Latex[length=2mm]}, thick, vertactu},
  fluxneg/.style={-{Latex[length=2mm]}, thick, rougeactu},
  etiq/.style={font=\scriptsize, inner sep=1pt}
}

% ------------------------------------------------------------------ titres
\usepackage{titlesec}
\titleformat{\section}{\Large\bfseries\sffamily\color{bleuactu}}{\thesection}{0.7em}{}
\titleformat{\subsection}{\large\bfseries\sffamily\color{bleuactu!80!black}}{\thesubsection}{0.6em}{}

\pagestyle{fancy}
\fancyhf{}
\fancyhead[L]{\small\sffamily\color{grisactu}Prérequis --- Comptabilité prudentielle}
\fancyhead[R]{\small\sffamily\color{grisactu}\thepage}
\renewcommand{\headrulewidth}{0.2pt}

\newcommand{\notation}[2]{\item[{$#1$}] #2}  % accolades : #1 peut contenir des crochets


% Version étudiant par défaut ; `make pdf` produit aussi td-corrige.pdf en définissant \Corrige.
\ifdefined\Corrige\correctiontrue\else\correctionfalse\fi

\fancyhead[L]{\small\sffamily\color{grisactu}Séance 2 --- Comptabilité prudentielle}
\newcommand{\meur}{\,M\euro}
\newcommand{\fx}[1]{\texttt{\small #1}}

\title{\vspace{-1.5cm}\sffamily\bfseries\color{bleuactu}
       TD 2 : contrôler les états de six assureurs\\[2mm]
       \large\mdseries\color{black}Avec un tableur, puis en Python, puis avec un agent d'IA
       \ifcorrection\\[1mm]\color{vertactu}Version corrigée\fi}
\author{\sffamily Fabrice Borel-Mathurin}
\date{\sffamily\small Master SAFIR --- 2026-2027}

\begin{document}
\maketitle
\thispagestyle{empty}

\begin{notion}{Le projet (2 heures, trois parties de 40 minutes)}
Vous êtes chargé de vérifier, avant leur remise au superviseur, quelques chiffres clés des états
de six assureurs fictifs : bilan, fonds propres par niveau, SCR, MCR et ratio de couverture.
Vous faites les contrôles une première fois sur tableur (partie A), puis vous les automatisez en
Python (partie B). Enfin, vous relisez un script produit par un assistant d'IA et vous rédigez la
consigne d'un agent chargé d'étendre vos contrôles (partie C).

\smallskip
\textbf{À rendre} : le classeur, le script Python, la consigne et une note de revue d'une page.
\end{notion}

\section{Données et règles}

\texttt{bilans-assureurs.csv} (séparateur « ; », montants en M\euro) donne, pour chaque entité :
l'actif, les provisions techniques, les dettes subordonnées et les autres passifs du bilan
prudentiel ; l'excédent d'actif sur passif déclaré ; les fonds propres de niveau 1, 2 et 3 ; le
SCR ; le MCR ; le ratio de couverture du SCR déclaré (en \%). \texttt{matrice-correlation.csv}
donne la matrice de corrélation entre modules du SCR de base (directive, annexe IV).

\smallskip
\noindent Règles à appliquer :
\begin{itemize}[nosep]
  \item fonds propres de base $=$ excédent d'actif sur passif $+$ dettes subordonnées ;
  \item couverture du SCR (règlement délégué, art.~82) : niveau 3 $< 15\,\%$ du SCR ; niveaux 2 $+$ 3
        $\leq 50\,\%$ du SCR ; niveau 1 $\geq 50\,\%$ du SCR ;
  \item couverture du MCR : niveau 2 $\leq 20\,\%$ du MCR ; niveau 3 exclu ; niveau 1 $\geq 80\,\%$ du MCR ;
  \item le MCR est compris entre 25\,\% et 45\,\% du SCR (directive, art.~129) ; le plancher absolu
        est supposé non contraignant ;
  \item un ratio déclaré est jugé faux s'il s'écarte de plus de 0,5 point du ratio recalculé.
\end{itemize}

\section{Partie A : le tableur}

\begin{exercice}[Un classeur qu'un autre peut auditer]
Importer le fichier (\emph{Données > À partir d'un fichier texte/CSV}, séparateur point-virgule)
dans un onglet \emph{Données} et le mettre sous forme de tableau. Créer trois autres onglets :
\emph{Calculs}, \emph{Contrôles}, \emph{Synthèse}. Aucune valeur saisie en dur dans une formule : les
paramètres (15\,\%, 50\,\%, 20\,\%, 0,5 point…) vont dans des cellules nommées.
\end{exercice}
\begin{corrige}
La séparation entrées / calculs / sorties et les paramètres nommés permettent de relire le classeur
et de changer une règle à un seul endroit. La mise sous forme de tableau donne des références
lisibles (\fx{[@scr]}) et étend les formules aux lignes ajoutées.
\end{corrige}

\begin{exercice}[Recalculer]
Pour chaque entité : l'excédent d'actif sur passif ; les fonds propres éligibles à la couverture
du SCR, puis le ratio ; les fonds propres éligibles à la couverture du MCR, puis le ratio.
\end{exercice}
\begin{corrige}
Formules (Excel en français, tableau \fx{T}, paramètres nommés) :
\begin{itemize}[nosep]
  \item niveau 3 éligible : \fx{=MIN([@tier3];Lim\_T3*[@scr])}
  \item niveau 2 éligible : \fx{=MIN([@tier2];Lim\_T2T3*[@scr]-[@t3\_elig])}
  \item ratio : \fx{=([@tier1]+[@t2\_elig]+[@t3\_elig])/[@scr]}
  \item éligibles MCR : \fx{=[@tier1]+MIN([@tier2];Lim\_T2\_MCR*[@mcr])}
\end{itemize}
\begin{center}
\small
\begin{tabular}{llrrrrr}
\toprule
 & Entité & Excédent & Éligibles SCR & Ratio SCR & Éligibles MCR & Ratio MCR \\
\midrule
A & Gamma   & 95  & 125 & 156,3\,\% & 91   & 303\,\% \\
B & Delta   & 65  & 85  & 100,0\,\% & 70   & 280\,\% \\
C & Epsilon & 80  & 59  & 98,3\,\%  & 50   & 250\,\% \\
D & Zêta    & 100 & 150 & 125,0\,\% & 105  & 420\,\% \\
E & Êta     & 25  & 35  & 87,5\,\%  & 22,4 & 187\,\% \\
F & Thêta   & 40  & --- & ---       & 40   & 400\,\% \\
\bottomrule
\end{tabular}
\end{center}
Pour C, le niveau 3 (30) est plafonné à 9 : le ratio passe de 133\,\% (déclaré, limites oubliées)
à 98\,\%, et le SCR n'est plus couvert.
\end{corrige}

\begin{exercice}[Contrôler]
Dans l'onglet \emph{Contrôles}, une colonne par contrôle, à valeur \fx{VRAI} en cas d'anomalie :
\begin{enumerate}[nosep]
  \item[C0] une donnée manque ;
  \item[C1] l'excédent déclaré diffère de l'excédent recalculé ;
  \item[C2] la somme des niveaux diffère des fonds propres de base ;
  \item[C3] le ratio déclaré est faux ;
  \item[C4] le MCR sort du corridor de 25 à 45\,\% du SCR ;
  \item[C5] le SCR n'est pas couvert ; C6 le MCR n'est pas couvert.
\end{enumerate}
Mettre en couleur les anomalies (mise en forme conditionnelle) et compter, dans la \emph{Synthèse},
les anomalies par entité (\fx{NB.SI}).
\end{exercice}
\begin{corrige}
\begin{center}
\small
\begin{tabular}{lll}
\toprule
Entité & Anomalies & Lecture \\
\midrule
A & aucune & --- \\
B & C1 & excédent déclaré 75, recalculé 65 \\
C & C3, C5 & ratio déclaré 133,3\,\% sans les limites ; 98,3\,\% avec \\
D & C4 & MCR $=$ 20,8\,\% du SCR \\
E & C5 & ratio 87,5\,\% : SCR non couvert (déclaré comme tel) \\
F & C0 & SCR et ratio manquants : C3 à C5 ne peuvent pas être évalués \\
\bottomrule
\end{tabular}
\end{center}
C5 pour E n'est pas une erreur de déclaration mais une alerte : l'entité doit présenter un plan
de rétablissement (directive, art.~138).
\end{corrige}

\begin{exercice}[Agréger des SCR]
Pour l'Assureur Gamma, les SCR des modules valent : marché 60, contrepartie 5, vie 30, santé 0,
non-vie 0 (M\euro). Calculer le SCR de base $\sqrt{s^\top C s}$ avec la matrice fournie, et le
comparer à la simple somme.
\end{exercice}
\begin{corrige}
\fx{=RACINE(PRODUITMAT(PRODUITMAT(TRANSPOSE(s);C);s))} avec \fx{s} le vecteur colonne nommé et
\fx{C} la matrice nommée : \textbf{75,17\meur}, contre 95 en simple somme ; la diversification vaut
19,83\meur. Le risque des immobilisations incorporelles s'ajouterait hors de la racine
(règlement délégué, art.~87).
\end{corrige}

\section{Partie B : Python}

\begin{exercice}[Automatiser]
Partir du notebook \emph{Python express} de la séance. Reproduire en pandas les calculs et les
contrôles de la partie A, sans boucle sur les lignes (\fx{np.minimum} applique une limite à toute
une colonne). Produire un tableau des anomalies par entité, et vérifier par des assertions que vous
retrouvez exactement les résultats du tableur.
\end{exercice}
\begin{corrige}
Script complet : \texttt{corriges/td-2\_controles.py}. Le cœur tient en trois lignes :

\smallskip
\fx{t3 = np.minimum(d["tier3"], 0.15 * d["scr"])}\\
\fx{t2 = np.minimum(d["tier2"], 0.50 * d["scr"] - t3)}\\
\fx{ratio = (d["tier1"] + t2 + t3) / d["scr"]}

\smallskip
Les assertions du script reprennent les résultats de la partie A (ratios de A, B, C ; anomalies
de chaque entité ; BSCR de 75,17) : si le code change, une régression se voit immédiatement.
\end{corrige}

\begin{exercice}[Les cas limites]
Que fait votre script de l'entité F ? D'un ratio d'exactement 100\,\% (entité B) ? Comment
comparez-vous un ratio recalculé à un ratio déclaré arrondi au dixième ? Justifier chaque choix.
\end{exercice}
\begin{corrige}
F doit rester dans le rapport, signalée « donnée manquante » : la supprimer ferait disparaître
l'entité la plus douteuse. Un ratio de 100\,\% couvre le SCR (fonds propres éligibles $\geq$ SCR) :
le test est \fx{>= 1}, pas \fx{> 1}. Deux nombres à virgule ne se comparent pas avec \fx{==} :
on compare l'écart à une tolérance (ici 0,5 point).
\end{corrige}

\section{Partie C : travailler avec un agent d'IA}

\begin{exercice}[Relire un script produit par une IA]
Le fichier \texttt{script-ia-a-relire.py} a été produit par un assistant d'IA à partir d'une
consigne d'une phrase. Relevez toutes ses erreurs (il y en a au moins huit). Pour chacune : sa
nature (réglementaire, données, programmation, communication) et son effet sur les résultats.
Exécutez-le pour le vérifier, après avoir corrigé ce qui l'empêche de tourner.
\end{exercice}
\begin{corrige}
\small
\begin{tabular}{p{0.3cm}p{6.6cm}p{7.2cm}}
\toprule
 & Erreur & Effet \\
\midrule
1 & « Article 82 de la directive » & La règle est à l'art.~82 du règlement délégué : référence fausse, énoncée avec assurance \\
2 & Chemin absolu \fx{C:/Users/…} & Le script ne tourne sur aucune autre machine \\
3 & \fx{dropna()} sans avertissement & F disparaît ; le rapport annonce « 5 entités analysées » \\
4 & Niveau 3 plafonné à 15\,\% des fonds propres, pas du SCR & C passe à 103,3\,\% : déclaré couvert à tort \\
5 & Niveau 2 plafonné seul à 50\,\% du SCR, pas niveaux 2 $+$ 3 & A passe à 162,5\,\% : son ratio déclaré, juste, est signalé faux \\
6 & Minimum de niveau 1 (50\,\% du SCR, 80\,\% du MCR) non vérifié & Sans effet sur ces données, faux dans d'autres cas \\
7 & MCR : niveau 3 compté, niveau 2 non plafonné & E : 292\,\% au lieu de 187\,\% \\
8 & Couverture testée par \fx{> 1} & B, à 100\,\%, déclaré non couvert \\
9 & Comparaison de flottants par \fx{==} & Fragile : un ratio déclaré arrondi sera signalé faux \\
10 & Excédent et corridor du MCR non contrôlés ; « toutes les anomalies ont été détectées » & B et D passent inaperçus ; la conclusion est fausse et rien ne la justifie \\
\bottomrule
\end{tabular}

\smallskip
\normalsize
La consigne d'une phrase ne donnait ni les règles, ni les cas limites, ni de quoi vérifier : l'agent
a comblé les trous avec des hypothèses plausibles et fausses.
\end{corrige}

\begin{exercice}[Rédiger la consigne d'un agent]
Rédiger la consigne à donner à un agent d'IA pour qu'il ajoute à \emph{votre} script les contrôles
C1 (excédent) et C4 (corridor du MCR). Elle précise : l'objectif, le contexte et les données, les
contraintes, les critères d'acceptation, ce que l'agent ne doit pas faire, ce qu'il doit livrer.
Si vous disposez d'un agent, lancez-le, puis rédigez une note de revue : ce que vous avez vérifié,
comment, ce que vous avez accepté ou refusé.
\end{exercice}
\begin{corrige}
\small
\textbf{Objectif} : ajouter à \texttt{controles.py} deux contrôles, C1 « excédent déclaré $\neq$
actif $-$ provisions techniques $-$ dettes subordonnées $-$ autres passifs » et C4 « MCR hors de
25--45\,\% du SCR » (directive, art.~129). \textbf{Contexte} : \texttt{bilans-assureurs.csv},
données fictives, M\euro, séparateur « ; ». \textbf{Contraintes} : pandas et NumPy seulement ;
ne pas modifier les contrôles existants ; ne supprimer aucune ligne (une donnée manquante est
signalée par C0) ; comparaisons avec tolérance. \textbf{Critères d'acceptation} : C1 signale B et
lui seul ; C4 signale D et lui seul ; F reste dans le rapport ; une assertion par critère ; le
script s'exécute depuis le dossier de la séance. \textbf{Ne pas} : citer un article sans l'avoir
vérifié dans le texte, conclure que « toutes les anomalies » sont détectées. \textbf{Livrer} :
les modifications (diff), la sortie du script, la liste des vérifications faites.
\end{corrige}

\begin{exercice}[Les règles du jeu en entreprise]
Dans un assureur, quelles données pouvez-vous confier à un agent d'IA ? Qui répond d'un chiffre
qu'il a calculé ? Que faut-il conserver pour qu'un auditeur ou le superviseur puisse refaire le
chemin ? Quels textes encadrent ces usages ?
\end{exercice}
\begin{corrige}
\textbf{Données} : uniquement ce que la politique de l'entreprise autorise pour l'outil utilisé ;
jamais de données personnelles hors d'un cadre conforme au RGPD (règlement (UE) 2016/679), ni
d'informations confidentielles dans un service non approuvé. \textbf{Responsabilité} : l'entreprise
et la personne qui valide restent responsables, comme pour toute activité externalisée (directive,
art.~49) ; le système de gouvernance doit garantir une gestion saine et prudente (art.~41).
\textbf{Traçabilité} : code et consignes versionnés, tests, revue par une seconde personne, sources
réglementaires vérifiées. \textbf{Textes} : règlement (UE) 2024/1689 sur l'IA (systèmes à haut
risque, dont l'évaluation des risques et la tarification en assurance vie et santé, annexe~III) ;
règlement (UE) 2022/2554 (DORA) pour le risque informatique et les prestataires tiers, applicable
depuis le 17 janvier 2025.
\end{corrige}

\begin{generalisation}{Ce que le TD simplifie}
Six entités et quelques colonnes au lieu des états complets ; des fonds propres de base sans
déduction ni élément restreint ; un plancher absolu du MCR ignoré. Les vrais contrôles de cohérence
des états sont les règles de validation publiées par l'EIOPA avec la taxonomie XBRL.
\end{generalisation}

\end{document}
